% Section: Hyperparameter Optimization for Non-Stationary Routing
% Include from paper/sections/results.tex (main paper, not appendix)

\subsection{Hyperparameter Optimization for Non-Stationary Routing}
\label{sec:hparam_sweep}

\paragraph{The challenge.}
BanditGPT's three core hyperparameters---exploration coefficient
$\alpha$, prior strength $\neff$, and forgetting factor
$\gamma$---interact non-trivially with the stationarity of the
deployment environment.
Standard cross-validation assumes i.i.d.\ data and optimises a
single stationary metric (e.g., reward or Pareto AUC), which
systematically selects $\gamma{=}1.0$ (infinite memory)
because forgetting \emph{always} hurts stationary performance.
However, $\gamma{=}1.0$ fails catastrophically under distribution
shift: as shown in Section~\ref{sec:nonstationarity}, the
no-forgetting condition accumulates $58.8\%$ more cumulative regret
than the tuned $\gamma{=}0.997$ after a reward swap.
Conversely, optimising solely for adaptability yields
$\gamma \ll 1$ with excessive forgetting that destabilises the
stationary policy.
This creates a fundamental tension absent from supervised routing:
\emph{the same hyperparameter must simultaneously protect stationary
efficiency and enable non-stationary recovery}.

Existing bandit-based LLM routers do not address this tension.
PILOT~\cite{panda2025pilot} grid-searches the exploration rate
$\alpha$ to maximise a single reward metric on held-out data,
with no forgetting mechanism and no non-stationarity consideration
in the selection protocol.
BaRP~\cite{wang2025barp} and LLM Bandit~\cite{li2025llmbandit}
omit forgetting entirely.
MixLLM~\cite{wang2025mixllm} delegates adaptation to periodic
neural-network retraining rather than tuning bandit-specific
parameters.
In the general bandit literature, the CDT framework~\cite{kang2024cdt}
proposes online hyperparameter tuning via a double-layer bandit that
treats $\alpha$ selection as a non-stationary continuum-armed problem,
extending earlier discrete methods~\cite{bouneffouf2020oplinucb}.
However, CDT remains \emph{single-objective} (maximise cumulative
reward) and does not tune the forgetting factor $\gamma$ or address
the efficiency--adaptability tradeoff.
The foundational Discounted UCB analysis by
Garivier and Moulines~\cite{garivier2011switching} derives $\gamma$
theoretically from the number of breakpoints, but this quantity is
unknown in every production deployment.
To the best of our knowledge, no prior work---in LLM routing or in
the broader non-stationary bandit literature---proposes a
multi-objective protocol for jointly selecting $(\alpha, \neff,
\gamma)$ to balance stationary efficiency against non-stationary
recovery.

\paragraph{Multi-objective formulation.}
We resolve this tension via an
$\epsilon$-constraint method~\cite{haimes1971bicriterion} that
treats hyperparameter selection as a bi-objective optimisation
problem over the grid
$\alpha \in \{0.01, 0.05, 0.1, 0.25, 0.5, 1.0\}$,
$\neff \in \{1, 10, 50, 200, 1000, 5000\}$,
$\gamma \in \{0.995, 0.997, 0.999, 1.0\}$:
%
\begin{enumerate}[nosep,leftmargin=*]
  \item \textbf{Primary objective} (efficiency):
    Maximise budget-paced Pareto AUC on the validation split, computed
    by sweeping seven log-spaced budget targets with the BudgetPacer
    active, building per-seed Pareto frontiers (with fixed-model
    endpoints), and averaging AUC across 10~seeds.

  \item \textbf{Constraint}:
    Retain only configurations whose budget-paced AUC is within
    $\epsilon{=}5\%$ of the grid maximum, forming a \emph{feasible set}
    that protects stationary routing quality.

  \item \textbf{Secondary objective} (adaptability):
    Among feasible configurations, minimise the average Phase-2
    cumulative regret under a synthetic reward/cost swap
    (identical to the non-stationary protocol of
    Section~\ref{sec:reward_shift}, averaged over all
    $\binom{K}{2}$ arm-pair swaps and 10~seeds).
\end{enumerate}
%
This formulation gives budget-pacing \emph{priority}---no
configuration that materially degrades stationary routing can be
selected---while using the secondary objective to break ties among
near-optimal configurations by their ability to recover from
distribution shift.
The $\epsilon$ parameter gives the practitioner an explicit,
auditable knob: it specifies the maximum stationary quality they are
willing to sacrifice in exchange for non-stationary robustness.

\paragraph{Results.}
Table~\ref{tab:hparam_selection} summarises the selection outcome
for both BanditGPT (warmup priors) and Tabula Rasa (cold start).

\begin{table}[t]
\centering
\caption{%
  \textbf{$\epsilon$-constraint hyperparameter selection.}
  AUC-only optimisation selects $\gamma{=}1.0$ (no forgetting) for
  both variants, maximising stationary AUC but doubling Phase-2
  regret under reward shift.
  The $\epsilon$-constraint method ($\epsilon{=}5\%$) trades
  ${<}0.4\%$ AUC for ${\sim}50\%$ lower adaptation regret by
  selecting moderate forgetting factors.
  Metrics: budget-paced Pareto AUC (val, 10~seeds) and Phase-2
  cumulative regret (val, averaged over all arm-pair swaps,
  10~seeds).%
}
\label{tab:hparam_selection}
\small
\begin{tabular}{@{}llcccccc@{}}
\toprule
\textbf{Variant} & \textbf{Method}
  & $\boldsymbol{\alpha}$ & $\boldsymbol{\neff}$
  & $\boldsymbol{\gamma}$
  & \textbf{BP AUC} & \textbf{P2 Regret} \\
\midrule
\multirow{2}{*}{BanditGPT}
  & AUC-only     & $0.01$ & $1000$ & $1.000$ & $0.929$ & $92.5 \pm 45.1$ \\
  & $\epsilon$-constr.\ & $0.10$ & $10$   & $0.997$ & $0.926$ & $\mathbf{46.3 \pm 6.0}$ \\[3pt]
\multirow{2}{*}{Tabula Rasa}
  & AUC-only     & $0.25$ & ---    & $1.000$ & $0.928$ & $93.1 \pm 39.4$ \\
  & $\epsilon$-constr.\ & $0.01$ & ---    & $0.999$ & $0.927$ & $\mathbf{46.7 \pm 40.5}$ \\
\bottomrule
\end{tabular}
\end{table}

AUC-only optimisation selects $\gamma{=}1.0$ for \emph{both}
variants---exactly the infinite-memory pathology that
Section~\ref{sec:nonstationarity} shows is catastrophic under
drift.
The $\epsilon$-constraint method instead selects $\gamma{=}0.997$
(BanditGPT, effective memory ${\sim}333$ steps) and
$\gamma{=}0.999$ (Tabula Rasa, ${\sim}1{,}000$ steps),
trading ${<}0.4\%$ budget-paced AUC for a ${\sim}50\%$ reduction
in Phase-2 regret ($92.5 \to 46.3$ and $93.1 \to 46.7$).
On the held-out test split, the selected configurations achieve
Pareto AUC of $0.924$ (BanditGPT) and $0.926$ (Tabula Rasa),
confirming that the val-selected hyperparameters generalise.

Three additional patterns merit attention:

\begin{enumerate}[nosep,leftmargin=*]
  \item \textbf{The variance gap.}
    BanditGPT's Phase-2 regret standard deviation under
    $\epsilon$-constraint selection is $6.0$---nearly $7\times$ lower
    than the AUC-only selection ($45.1$) and the Tabula Rasa
    $\epsilon$-constraint ($40.5$).
    Warmup priors anchor the initial policy, ensuring that the
    forgetting factor operates on a well-calibrated baseline rather
    than a noisy cold-start estimate.
    This reliability advantage is critical for single-deployment
    production use, where the operator does not average over seeds.

  \item \textbf{Prior strength shrinks.}
    AUC-only selects $\neff{=}1{,}000$ (strong priors), while
    $\epsilon$-constraint selects $\neff{=}10$ (light priors).
    Strong priors increase effective memory and
    resist forgetting---beneficial for stationary AUC but
    counterproductive when the environment shifts and stale priors
    must be discounted.
    The $\epsilon$-constraint method discovers this tradeoff
    automatically.

  \item \textbf{Exploration increases.}
    AUC-only selects $\alpha{=}0.01$ (minimal exploration),
    while $\epsilon$-constraint selects $\alpha{=}0.1$
    ($10\times$ higher).
    Under non-stationarity, the router must re-explore arms whose
    quality may have changed; a higher exploration coefficient
    accelerates this re-discovery.
    This is consistent with findings in the online hyperparameter
    tuning literature~\cite{kang2024cdt}, which observes that the
    optimal $\alpha$ is non-stationary itself.
\end{enumerate}

\paragraph{Novelty and significance.}
This $\epsilon$-constraint protocol is, to our knowledge, the first
multi-objective hyperparameter selection framework for non-stationary
bandit routing.
It addresses a gap that spans both the LLM routing and the general
bandit literatures:
existing systems either assume stationarity in their tuning
protocol~\cite{panda2025pilot,wang2025barp,li2025llmbandit},
delegate adaptation to model retraining rather than tuning bandit
parameters~\cite{wang2025mixllm}, or derive $\gamma$ from
theoretical quantities (number of breakpoints) that are unknown in
production~\cite{garivier2011switching}.
The CDT framework~\cite{kang2024cdt} demonstrates that online
hyperparameter tuning is feasible for contextual bandits, but it
optimises a single reward objective and does not address the
stability--plasticity tradeoff that arises when the router must
simultaneously serve stationary traffic well and recover from shifts.

Our contribution is complementary: rather than tuning $\alpha$
online during execution (as CDT does), we propose a
\emph{deployment-time} protocol that jointly selects
$(\alpha, \neff, \gamma)$ using a multi-objective criterion
that explicitly encodes the dual requirements of
production routing---efficiency under stationarity and resilience
under drift.
The empirical result that single-objective tuning
\emph{systematically} selects $\gamma{=}1.0$ and that this doubles
Phase-2 regret is itself a cautionary finding: any practitioner
deploying a non-stationary bandit router who tunes via standard
cross-validation will converge on the worst possible forgetting
factor.

\paragraph{Practical implications.}
From a deployment perspective, the protocol eliminates the need to
hand-tune forgetting factors on live traffic---a practice that risks
degrading routing quality during the tuning period and offers no
reproducibility.
The practitioner specifies two quantities: (i)~a grid of candidate
hyperparameters, and (ii)~the tolerance $\epsilon$ (how much
stationary quality they are willing to sacrifice for robustness).
The output is a single configuration that is
near-optimal for budget-paced routing and maximally adaptive within
that quality budget.
Because the secondary objective uses a synthetic reward swap
(Section~\ref{sec:reward_shift}), the protocol requires no live
non-stationary data and can be run entirely on a static validation
split before deployment.
This makes the tuning procedure safe, offline, and repeatable---a
necessary property for production systems subject to change-management
governance.
